% Nouveau EGA V (cote 25): archive cote 25, site page 85 (PDF page 86).
% typescript. Draft compared: none.
% AI-transcribed by Claude Opus 5.5 (claude-opus-5-5, adaptive thinking, effort high), Message Batches API, batch ega25b (msgbatch_01A7n5YuXubUPNK1zsepi8RG), 2026-11-13; not yet checked by eye.
\scanpage{25/85}
\mspage{79}

n'est autre que celui qui associe à toute section de  P - C  l'unique
sous-variété linéaire  L de P, de codi\unsure{mmm}sion m, contenant à la fois C
et ladite section  (bien entendu, par ``contenir la section'' on entend\xout{x} qu\unsure{e}
% NOTE: the end of this line is cut by the right edge of the scan; only "qu" and part of the next letter are visible
la section se factorise par L).

Si maintenant on a un f: $X \longrightarrow P$,  il y a lieu de considérer le
composé
\[ X - f^{-1}(C) \longrightarrow\ P - C \ \longrightarrow\ Q(C) \]
qu'on pourra appeler projection conique de X, relativement à f et au cen\xout{x:}%
tre C\hand{,} \inserted{notée \unsure{$p^{f}_{C}$|$p^{\Sigma}$} ou $p^{X}_{C}$ ou simplement $p_{C}$.} \hand{(}On fera attention qu'elle n'est pas en général définie sur X tout
entier, d e façon précise elle l'est sss $f^{-1}(C)=\emptyset$ i.e. f(X) ne rencont re
pas le centre de projection C. \hand{\unsure{)}}

Nous allons donner une autre interprétation de ce morphisme, en
termes de constructions utilisées dans les n$^{\text{o}}$s antérieurs.  Pour ceci,
avec les notations introduites par ailleurs, considérons
\diagram{
\begin{tikzcd}
X & \underline{X}^{(m)}_{\xout{\unsure{Q}}Q(C)} = \underline{X}^{(m)}x_{\mathrm{Grass}^{m}}Q(C) \arrow[l, "q"'] \arrow[d, "p\xout{\unsure{z}}"'] \\
 & Q(C) \quad .
\end{tikzcd}
}
Notons d'autre part que  q  induit un  isomorphisme
\[ q':q^{-1}(X-f^{-1}(C)) \xrightarrow{\ \sim\ } X - f^{-1}(C) \ , \]
et il est immédiat que  $p_{C}$n'est autre  que   $p'\ q'^{-1}$ , où p' est la
restriction de  p à  $q^{-1}(X-f^{-1}(\xout{\unsure{Q}}C))$.  On peut dire encore, utilisant q'
pour faire une identificatio\xout{\unsure{n}}n pure et simp\unsure{be|le}, que  $p_{C}$ est la restriction
du morphisme p  à  $X - f^{-1}(C) \ \subset\ \underline{X}^{(m)}_{Q(C)}$ . Pour cette raison, il est
commode de dénoter encore par  $p^{X}_{C}$ \inserted{ou $p_{C}$}, et d'appeler  (mettons)  \uline{projection}
\uline{conique}  \uline{étendue}  de X\hand{,} relativement à $f:X\rightarrow P$, de centre C, le morphisme

\mspage{85}
